\section*{अध्याय \ref{ch:b3:structural-analysis} --- संरचनात्मक विश्लेषण}

\begin{solution}{exo:b3:structural-analysis:1}
\omterm{def:b3:structural-analysis:stereo}{ध्रुव} का दिगंश $40 + 180 = \ang{220}$ है और नमन $90 - 30 = \ang{60}$ है।
तल की नति $90 - 20 = \ang{70}$ है और नति-दिशा $315 - 180 = \ang{135}$ है: $135/70$।
\end{solution}

\begin{solution}{exo:b3:structural-analysis:2}
$\theta = \ang{60}$। सम-कोणीय: $10\tan\ang{30} = \qty{5.77}{cm}$;
सम-क्षेत्रफल: $10\sqrt2\sin\ang{30} = \qty{7.07}{cm}$।
\end{solution}

\begin{solution}{exo:b3:structural-analysis:3}
$\vec n_1 = (\sin\ang{30}, 0, \cos\ang{30})$, $\vec n_2 = (-\sin\ang{30}, 0,
\cos\ang{30})$, $\vec n_1\cdot\vec n_2 = -0.25 + 0.75 = 0.5$: दोनों \omterm{def:b3:structural-analysis:stereo}{ध्रुव}
\ang{60} दूर हैं, अंतर-पार्श्व कोण \ang{120} है, अतः यह एक खुला वलन है।
\end{solution}

\begin{solution}{exo:b3:structural-analysis:4}
$\vec d = (\cos\ang{60}, 0, -\sin\ang{60}) = (0.5, 0, -0.866)$ और $\vec u =
0.707\,\vec s + 0.707\,\vec d = (0.354, 0.707, -0.612)$: दिगंश
$\operatorname{atan2}(0.354, 0.707) = \ang{26.6}$, नमन $\arcsin0.612 =
\ang{37.8}$। रेक \ang{-90} (सामान्य) और 0 (बाएँ-पार्श्वीय) के बीच है: तिर्यक
सामान्य भ्रंश, जिसमें बाएँ-पार्श्वीय घटक है।
\end{solution}

\begin{solution}{exo:b3:structural-analysis:5}
\begin{gather*}
  \vec n_1 = (0.5567, 0.3214, 0.7660),\qquad \vec n_2 = (-0.4698, -0.1710, 0.8660),\\
  \vec n_1\times\vec n_2 = (0.4093, -0.8420, 0.0558),
\end{gather*}
जिसकी लंबाई $0.9379$ है। इसका विपरीत, $-\vec n_1\times\vec n_2$, नीचे की ओर इंगित करता है:
दिगंश \ang{334.1}, नमन $\arcsin(0.0558/0.9379) = \ang{3.4}$। $\vec n_1\cdot\vec
n_2 = 0.3469$, अभिलंबों के बीच \ang{69.7}, अंतर-पार्श्व कोण \ang{110.3}:
खुला वलन।
\end{solution}

\begin{solution}{exo:b3:structural-analysis:6}
नमन \ang{60} से अधिक का अर्थ है $\theta < \ang{30}$। सम-क्षेत्रफल: $r^2 =
2\sin^2\ang{15} = 0.134$, \qty{13.4}{\percent} जाल का; सम-कोणीय:
$\tan^2\ang{15} = 0.072$, \qty{7.2}{\percent}। गोले पर टोपी
$1 - \cos\ang{30} = 0.134$ भाग गोलार्ध का घेरती है: केवल सम-क्षेत्रफल जाल
उसे उसके सही अनुपात में दिखाता है, इसलिए \omterm{def:b3:structural-analysis:stereo}{ध्रुव} उसी पर गिने जाते हैं।
\end{solution}

\begin{solution}{exo:b3:structural-analysis:7}
\begin{gather*}
  T_{xx} = \frac{2\sin^2\ang{20}}3 = 0.078,\qquad
  T_{zz} = \frac{2\cos^2\ang{20} + 1}3 = 0.922,\\
  T_{xz} = -\frac{\sin\ang{20}\cos\ang{20} - \sin\ang{20}\cos\ang{20}}3 = 0,
\end{gather*}
और $y$ वाला हर पद 0 है। आइगेनमान $0.922$ (ऊर्ध्वाधर),
$0.078$ (पूर्व) और $0$ (उत्तर): \omterm{def:b3:structural-analysis:pi}{बीटा-अक्ष} क्षैतिज है, उत्तर--दक्षिण।
तीनों \omterm{def:b3:structural-analysis:stereo}{ध्रुव} एक ही तल में हैं, इसलिए किसी का भी अपने अभिलंब के अनुदिश घटक नहीं है,
और $\vec u^{\mathsf T}T\vec u = 0$ है, जहाँ $\vec u$ उसके अनुदिश है।
\end{solution}

\begin{solution}{exo:b3:structural-analysis:8}
$a = 2$, $b = 1.5$, $k = 1/0.5 = 2$: संकुचन। $Z = 1/(1.8\times1.1) =
0.505$, $a = 1.636$, $b = 2.178$, $k = 0.636/1.178 = 0.54$: चपटापन।
\end{solution}

\begin{solution}{exo:b3:structural-analysis:9}
$S_1/S_2 = 2.25$ और $S_1S_2 = 1$: $S_1 = 1.5$, $S_2 = 0.667$; लंबी अक्ष
\qty{50}{\percent} बढ़ाई गई थी। रिक्ति के अर्ध-अक्ष $1.5\times0.4 =
\qty{0.60}{mm}$ और $0.667\times0.4 = \qty{0.27}{mm}$ हैं।
\end{solution}

\begin{solution}{exo:b3:structural-analysis:10}
सम-कोणीय: $1/(2\cos^2\ang{30}) = 0.667$ त्रिज्या के अनुदिश और
$\tan\ang{30}/\sin\ang{60} = 0.577/0.866 = 0.667$ चारों ओर: बराबर, समकोणिक।
सम-क्षेत्रफल: $\mathrm dr/(R\,\mathrm d\theta) = (\sqrt2/2)\cos\ang{30} = 0.612$
और $\sqrt2\sin\ang{30}/\sin\ang{60} = 0.707/0.866 = 0.816$; गुणनफल
$0.500$ है, जो प्रमेय का क्षेत्रफल गुणक है। एक लघु वृत्त ऐसे अक्षों वाला अंडाकार बन जाता है,
जिनका अनुपात $0.816/0.612 = 1.33$ है, और लंबा अक्ष केंद्र के चारों ओर है।
\end{solution}

\begin{solution}{exo:b3:structural-analysis:11}
भ्रंश ऊर्ध्वाधर है, इसलिए उसका अभिलंब क्षैतिज है और
ऊर्ध्वाधर $\sigma_2$ है; प्रमेय के भाग (ख) से सर्पण इसके लंबवत होता है,
$\sigma_2$ के लंबवत, क्षैतिज है, हर $\Phi$ के लिए। $\sigma_1 = 1$ उत्तर के अनुदिश होने पर,
$\sigma_3 = 0$ पूर्व के अनुदिश: $S\vec n = (0, -0.342, 0)$, $\vec n\cdot S\vec n =
0.117$, $\vec\tau = (0, -0.342, 0) - 0.117\,(0.940, -0.342, 0) = (-0.110,
-0.302, 0)$ और $-\vec\tau = (0.110, 0.302, 0)$, \ang{020} की ओर। खंड,
जो $\vec n$ की ओर है, यानी पूर्वी पक्ष, स्ट्राइक के अनुदिश उत्तर-उत्तर-पूर्व चलता है:
पश्चिमी पक्ष से देखने पर वह बाईं ओर चलता है, यह एक बाएँ-पार्श्वीय भ्रंश है (रेक
0)। $\sigma_2$ ने गणना में कोई भूमिका नहीं निभाई।
\end{solution}

\begin{solution}{exo:b3:structural-analysis:12}
$S = (7.5 - 6.0)/1.0 = \qty{1.5}{km}$; $h = 6.0/1.5 = \qty{4.0}{km}$ निचले
क्षितिज से नीचे। अनुभाग की लंबाई $20 + 1.5 = \qty{21.5}{km}$ थी, और उसमें परिवर्तन
$-1.5/21.5 = \qty{-7.0}{\percent}$ रहा।
\end{solution}

\begin{solution}{pb:b3:structural-analysis:1}
\textbf{1.} \omterm{def:b3:structural-analysis:stereo}{ध्रुव} का दिगंश \ang{255} है और नमन \ang{77} है; $\theta =
\ang{13}$ और $r = 10\sqrt2\sin\ang{6.5} = \qty{1.60}{cm}$, केंद्र से,
\ang{255} की ओर।

\textbf{2.} $\vec n = (\sin\ang{13}\sin\ang{75}, \sin\ang{13}\cos\ang{75},
\cos\ang{13}) = (0.2173, 0.0582, 0.9744)$। लगभग समांतर छह इकाई सदिश
जुड़कर लगभग 6 देते हैं: $|\vec S_E| = 5.9974$।

\textbf{3.} पूर्वी पार्श्व: $\bar\alpha = \operatorname{atan2}(1.2583, 0.2017) =
\ang{80.9}$ और $\bar\delta = \arccos(5.8604/5.9974) = \ang{12.3}$; $\bar R
= 5.9974/6 = 0.99956$ और प्रकीर्णन $\sqrt{2\times0.00044} =
\qty{0.0297}{rad}$, या \ang{1.7}। पश्चिमी पार्श्व: $\bar\alpha =
\operatorname{atan2}(-1.2445, 0.1886) = \ang{278.6}$ और $\bar\delta =
\arccos(5.8636/5.9972) = \ang{12.1}$; $\bar R = 0.99953$, प्रकीर्णन \ang{1.8}।

\textbf{4.}
\[
  \vec S_E\times\vec S_W = (0.0774, -14.6714, 0.4883),
\]
 लंबाई $14.680$ है। इसका विपरीत सदिश नीचे की ओर इंगित करता है: दिगंश $\operatorname{atan2}(-0.0774, 14.6714) = \ang{359.7}$ और
 नमन $\arcsin(0.4883/14.680) = \ang{1.9}$।

\textbf{5.} माध्य अभिलंबों के बीच के कोण की कोज्या
$32.835/35.968 = 0.9129$ है, अर्थात योगों के बिंदु-गुणनफल को उनकी
लंबाइयों से भाग देने पर: अभिलंबों के बीच \ang{24.1}, अंतर-पार्श्व कोण \ang{155.9}:
यह एक खुला वलन है, जो लगभग मंद वलन के निकट है।

\textbf{6.}
\begin{gather*}
  \vec m_E = (0.2098, 0.0336, 0.9772),\qquad \vec m_W = (-0.2075, 0.0314, 0.9777),\\
  \vec m_E - \vec m_W = (0.4173, 0.0022, -0.0005):
\end{gather*}
एक क्षैतिज अभिलंब \ang{089.7} की ओर है। अक्षीय तल ऊर्ध्वाधर है, जिसका स्ट्राइक
\ang{359.7} है, लगभग उत्तर--दक्षिण: वलन सीधा है।

\textbf{7.} मेखला वह तल है जिसका \omterm{def:b3:structural-analysis:stereo}{ध्रुव} \omterm{def:b3:structural-analysis:pi}{बीटा-अक्ष} है: नति-दिशा
\ang{179.7}, नति $90 - 1.9 = \ang{88.1}$, लगभग पूर्व--पश्चिम ऊर्ध्वाधर
तल। प्रत्येक पार्श्व में वलन-अक्ष होता है, अतः प्रत्येक अभिलंब उसके
लंबवत होता है और उसके लंबवत तल में स्थित है।

\textbf{8.} $0.04400 + 0.00144 + 0.95456 = 1.00000$। $\vec b^{\mathsf T}T\vec
b = 0.00038$, सबसे छोटा आइगेनमान: $\vec b$ पूर्णांकन की सीमा तक
आइगेनसदिश $\vec e_3$ है, और दो माध्य ध्रुवों से होकर जाने वाली मेखला
बारहों में से सर्वोत्तम-फ़िट मेखला है।

\textbf{9.} $\arcsin\sqrt{0.00038} = \arcsin0.0195 = \ang{1.1}$।

\textbf{10.} मेखला के अनुदिश \omterm{def:b3:structural-analysis:stereo}{ध्रुव} ऊर्ध्वाधर से लगभग \ang{12.2} पर स्थित हैं,
दोनों ओर; $\sin^2\ang{12.2} = 0.045$, उनके पूर्व--पश्चिम घटकों का
माध्य वर्ग, जो $\lambda_2$ है।

\textbf{11.} हाँ: \omterm{def:b3:structural-analysis:stereo}{ध्रुव} औसतन \ang{1.1} से मेखला से विचलित होते हैं, जो
एक पार्श्व के अभिलेखों के अपने माध्य के परितः \ang{1.7} प्रकीर्णन से अधिक नहीं है।

\textbf{12.} $\ang{1.1}/\sqrt{12} = \ang{0.3}$: नमन \ang{1.9} $\pm$
\ang{0.3} है।

\textbf{13.} \ang{-90} का रेक: शुद्ध सामान्य भ्रंश, नति-सर्पण। एंडरसन:
$\sigma_1$ ऊर्ध्वाधर, $\sigma_3$ नति-दिशा के अनुदिश क्षैतिज,
\ang{120}--\ang{300}, $\sigma_2$ स्ट्राइक के अनुदिश क्षैतिज,
\ang{030}--\ang{210}।

\textbf{14.} $120/60$ का अभिलंब उस ऊर्ध्वाधर तल में है, जिसका दिगंश
\ang{120} है, जिसमें $\vec e_1$ और $\vec e_3$ हैं: $n_2 = 0$। भाग (ख) के अनुसार
सर्पण $\sigma_2$ के लंबवत, अर्थात स्ट्राइक के, हर $\Phi$ के लिए है: यह
नति-सर्पण है। यह भ्रंश $\Phi$ के बारे में कुछ नहीं बताता; सूत्र में $n_2$ और
$\vec e_2\cdot\vec w$ दोनों शून्य हो जाते हैं, अतः अंश और हर दोनों लुप्त हो जाते हैं।

\textbf{15.} $\vec n = (\sin\ang{55}\sin\ang{150}, \sin\ang{55}\cos\ang{150},
\cos\ang{55}) = (0.4096, -0.7094, 0.5736)$; $n_1 = 0.574$; $n_2 = 0.5\times
0.4096 + 0.866\times(-0.7094) = -0.410$।

\textbf{16.} $\vec e_1\cdot\vec w = 0.198$ और $\vec e_2\cdot\vec w = 0.5\times
0.771 + 0.866\times0.605 = 0.909$, अतः $\Phi = -(0.574\times0.198)/(-0.410\times
0.909) = 0.1137/0.3727 = 0.305$।

\textbf{17.} $(0.305 + 0.294 + 0.301 + 0.306)/4 = 0.30$। रेक
एक अंश तक पूर्णांकित किए गए थे; इन भ्रंशों पर रेक का एक अंश $\Phi$ को लगभग 0.02
बदल देता है, अतः आधे अंश की पूर्णांकन-त्रुटि प्रसार को समझा देती है।

\textbf{18.} $\sigma_2 - \sigma_3 = 0.3(\sigma_1 - \sigma_3)$: $\sigma_2$
$\sigma_3$ के अधिक निकट है, $\sigma_1$ की तुलना में। दोनों क्षैतिज प्रतिबल
ऊर्ध्वाधर प्रतिबल से कहीं कम हैं; क्षेत्र मुख्यतः \ang{120} की ओर फैला था,
थोड़ा \ang{030} के अनुदिश।

\textbf{19.} अभिलंब $\vec n = (0.750, -0.433, 0.500)$ और सर्पण $\vec u =
(0.433, -0.250, -0.866)$, $(\vec n - \vec u)/\sqrt2$ का \ang{300} की ओर नमन
\ang{75} है: P; $(\vec n + \vec u)/\sqrt2$ का \ang{120} की ओर नमन \ang{15} है:
T। ये उसी ऊर्ध्वाधर तल में हैं जिसमें $\sigma_1$ और $\sigma_3$ हैं, परंतु
उनसे \ang{15} दूर, क्योंकि भ्रंश \ang{45} नहीं, \ang{30} का कोण
$\sigma_1$ के साथ बनाता है।

\textbf{20.} $\tan a = \tan\ang{12.3}\times\cos(\ang{90} - \ang{80.9}) =
0.2180\times0.9874 = 0.2153$: $a = \ang{12.1}$।

\textbf{21.} $L_0 = 8000/\cos\ang{12} = \qty{8178.7}{m}$, $S =
\qty{178.7}{m}$, सापेक्ष परिवर्तन $\cos\ang{12} - 1 =
\qty{-2.2}{\percent}$ है।

\textbf{22.} $4000\tan\ang{12} = \qty{850.2}{m}$; $A = \tfrac12\times8000\times
850.2 = \qty{3.40e6}{m^2}$।

\textbf{23.} $h = \num{3.40e6}/178.7 = \qty{19.0}{km}$। आवरण के आधार पर विलगन
होने पर $h$ कुछ सौ मीटर का और अतिरिक्त क्षेत्रफल
लगभग पचास गुना छोटा होता: अतः वलन वहाँ विलगित नहीं हैं।
या तो आधारभूत चट्टानें आवरण के साथ वलित हुईं, या अभिनतियाँ मूल
स्तर से नीचे धँस गईं, जिससे $A$ छोटा हो जाता है; अकेली काट यह नहीं बता सकती।

\textbf{24.} क्षैतिज विस्थापन $200/\tan\ang{60} = \qty{115.5}{m}$, सर्पण $200/\sin\ang{60}
= \qty{230.9}{m}$।

\textbf{25.} $\tan a = \tan\ang{60}\cos\ang{30} = 1.500$: $a = \ang{56.3}$, और
आभासी क्षैतिज विस्थापन $200/1.5 = \qty{133.3}{m}$। भ्रंश का सर्पण
नति के अनुदिश नीचे, \ang{120} की ओर है: \ang{120} के अनुदिश काट वास्तविक
लंबाइयाँ दिखाती है और वही पुनर्स्थापित की जानी चाहिए।

\textbf{26.} भ्रंश वलनों को काटता है, अतः वह युवा है और सबसे पहले पुनर्स्थापित किया जाता है,
ऊपरी भित्ति को भ्रंश के अनुदिश ऊपर की ओर \qty{230.9}{m} पीछे ले जाकर; वलन
फिर खोले जाते हैं। \omterm{def:b3:structural-analysis:restoration}{पिन रेखाएँ} अभिनतियों की अक्षीय सतहों पर रखी जाती हैं, निचली भित्ति में,
जहाँ सममित वलनों में संस्तरों के बीच कोई सर्पण नहीं था।

\textbf{27.} वार्डन पहाड़ियों का वलन-अक्ष
\textbf{\ang[text-series-to-math=true]{1.9}} नमन करता है, उत्तर की ओर (दिगंश
\ang{359.7}): क्षेत्र में निहित \ang{2} को
बारह अभिलेख जितनी अनुमति देते हैं, उस \ang{0.3} तक पुनः प्राप्त किया गया है।
\end{solution}
